% draft3.tex — ALIFE 2026 submission
% Compile from this directory with alifeconf.sty on the LaTeX path
% (copy alifeconf.sty from the parent folder, or set TEXINPUTS=../:)
% pdflatex draft3 && bibtex draft3 && pdflatex draft3 && pdflatex draft3

\documentclass[letterpaper]{article}

\usepackage{natbib,alifeconf}  %% Order is important
\usepackage{amsmath}
\usepackage{url,hyperref,cleveref}
\usepackage{booktabs}
\usepackage{graphicx}

\title{Allostatic Neural Substrate: Emergent Competence and\\
Native Interpretability Through Metabolically Constrained Local Agents}

\author{
    Daniel Pace$^{1}$ \\
    \mbox{}\\
    $^1$Independent Researcher \\
    daniel.a.pace725@gmail.com
}

\begin{document}

\maketitle

\begin{abstract}
We present REAL Neural Substrate, a learning architecture in which each network
node is an autonomous agent operating under a metabolic energy budget rather than
a passive unit updated by global gradient descent. Persistent routing substrate
accumulates through allostatic self-regulation and stigmergic substrate writing,
producing cross-task transfer, latent context inference, and difficulty-correlated
topology growth as emergent properties rather than engineered objectives. On a
real-world room occupancy dataset, the system reaches F1 of 0.952 against an MLP
baseline of 0.963, with substrate that transfers to a fresh system at 99.15\%
efficiency --- confirming that the routing substrate, not runtime state, is the
actual knowledge carrier.

We distinguish two roles that native interpretability plays in this architecture:
interpretability as structural \textbf{capability} --- local decisions, memory
structures, and metabolic states are directly readable because they are explicit
properties of local agents, not distributed consequences of global optimization
--- and interpretability as \textbf{development workflow} --- this readability was
the primary tool by which failure modes were diagnosed and architectural changes
were motivated throughout the system's development. We argue that building this
system, rather than formalizing it, revealed several things about learning and
agency that prior theory did not anticipate.
\end{abstract}

\section{Introduction}

The dominant paradigm in machine learning optimizes a global loss function by
propagating error gradients through every parameter simultaneously. This is
powerful, but it carries two structural costs. First, learned representations are
distributed across the full weight tensor: no individual component has a local
explanation for its behavior, because behavior is an emergent consequence of
millions of interacting parameters shaped by a non-local process. Second,
interpretability is retrospective. When the system fails, the failure appears in
aggregate output; understanding the mechanism requires external probe methods that
approximate explanations rather than read them.

Biological learning systems are organized differently. Neurons maintain local
energy budgets; synaptic change is local and conditional on pre- and post-synaptic
activity; learning accumulates in substrate structure that persists beyond any
episode. Allostasis --- regulation toward viable internal operating ranges --- is
the primary organizing principle. The substrate encoding a biological system's
history is, in principle, readable: its state at any moment reflects the causal
record of the organism's experience.

REAL (Relationally Embedded Allostatic Learning) Neural Substrate takes these
principles as architectural constraints rather than metaphors. The system was
initially developed on consumer hardware using real CPU cycles as the metabolic
currency, establishing metabolic constraint as a physical reality rather than an
analogy. In Phase 8, this grounding was extended to a multi-agent network: each
node is a local allostatic agent that observes only its immediate neighborhood,
selects actions from an explicit vocabulary under metabolic pressure, evaluates
its own behavior across six relational coherence dimensions, and consolidates
useful patterns into a persistent substrate that physically reduces the future
cost of behaviors that have earned feedback. No node has a global view. No global
gradient is computed.

This paper makes three connected claims. \textbf{Competence:} a network of such
agents, without any global objective, develops cross-task transfer, latent context
inference, difficulty-correlated topology growth, and near-MLP performance on a
real-world dataset it was not designed for. \textbf{Interpretability as
capability:} the architecture is natively inspectable because computation emerges
from agents with explicit local state, not from global function composition.
\textbf{Interpretability as development workflow:} this readability was the
primary development methodology --- and the experimental record documents it.

\section{Background}

\subsection{Biologically-Grounded Learning}

A persistent thread in ALife and cognitive science holds that biological learning
is fundamentally allostatic: organisms regulate toward viable internal operating
ranges, and behavioral adaptation is a consequence of that regulation rather than
a separately optimized objective. This contrasts with supervised learning, where
an external loss signal drives parameter update through a mechanism with no
plausible biological analogue.

Predictive processing frameworks \citep{rao1999predictive} propose that biological
neural circuits are fundamentally anticipatory, maintaining generative models of
expected inputs and updating on prediction error. REAL operationalizes this at the
node level: each agent runs a local recognize--predict--select loop in which
prediction error modulates the current selection context.

Stigmergy --- learning through modification of a shared environment rather than
through direct communication or centralized update --- underlies distributed
cognition in biological and artificial swarm systems
\citep{theraulaz1999stigmergy}. REAL's substrate is stigmergic in this sense: it
is written by action, read by future action, and persists across sessions.

Prior work on local learning rules has proposed alternatives to global credit
assignment, including contrastive Hebbian learning \citep{movellan1991contrastive}
and equilibrium propagation \citep{scellier2017equilibrium}. REAL differs from
this tradition in not attempting to approximate backpropagation with a local rule,
but in replacing the optimization objective entirely with allostatic coherence.

\subsection{Interpretability in Neural Systems}

The interpretability literature broadly divides into post-hoc methods and
transparent-by-design approaches. Post-hoc methods such as SHAP
\citep{lundberg2017shap}, LIME \citep{ribeiro2016lime}, and integrated gradients
\citep{sundararajan2017integrated} approximate explanations of trained models
through input perturbation and activation analysis. They are correlational: they
identify input sensitivity, not the decision procedure executed.

Transparent-by-design approaches --- decision trees, linear models, rule-based
systems --- achieve legibility at the cost of expressivity. REAL occupies a
distinct position: expressive enough to reach near-MLP performance, while
remaining natively interpretable as a structural consequence of the architecture,
not as a post-hoc addition.

\section{Architecture}

\subsection{The REAL Node}

Each node in REAL Neural Substrate runs an independent local loop: perceive,
select, execute, evaluate, consolidate. Three architectural properties are
central.

\textbf{Observation is strictly local.} A node sees its own ATP balance, the
packets currently present, the current context signal, and the substrate-key
indices of its own prior routing history. It has no access to any other node's
state or to a global objective.

\textbf{Selection is substrate-biased.} The base ATP cost of action $a$ under
context $k$ is discounted by accumulated substrate support:
\begin{equation}
  c_t(a \mid k) \;=\; c_0(a)\,\bigl(1 - \alpha \cdot M_s(a,\, k)\bigr)
  \label{eq:cost}
\end{equation}
where $c_0(a)$ is the base cost, $M_s(a,k) \in [0,1]$ is the substrate support
for action $a$ under context $k$, and $\alpha$ is the discount rate. Actions that
have historically earned feedback cost less to repeat --- allostasis in the
literal sense.

\textbf{Consolidation is explicit and inspectable.} Useful patterns migrate from a
fast episodic trace through a consolidated memory layer into the maintained
substrate $M_s$. The substrate is not a weight matrix; it is a set of named
entries that can be enumerated, compared across sessions, and selectively included
or excluded. A node's learning history is readable as structure.

\subsection{Coherence Evaluation}

Rather than measuring performance against an external loss, each node scores its
own cycle outcome across six relational dimensions:
\begin{equation}
  \Phi \;=\; \sum_{i=1}^{6} w_i\, \phi_i
  \label{eq:coherence}
\end{equation}
The six dimensions are: \textit{continuity} (identity persistence through action
sequences); \textit{vitality} (productive energy expenditure, penalizing both idle
and exhausted states); \textit{contextual fit} (alignment between action and
context signal); \textit{differentiation} (contrast with neighboring nodes,
penalizing redundancy); \textit{accountability} (coherence between intent and
observed outcome); and \textit{reflexivity} (behavioral revision following
coherence dips). Together, these constitute endogenous evaluation: the node asks
``am I maintaining coherent operation?'' rather than ``did an external observer
approve my output?''

\subsection{Phase 8: A Multi-Agent Routing Substrate}

In Phase 8, many REAL nodes operate simultaneously in a routing graph. Packets
enter at a source, must reach a sink with a context-dependent payload transform
applied, and the sink returns feedback upstream through the successful path. Nodes
on the path earn ATP; nodes that waste budget on fruitless actions are eventually
starved.

The durable learned structure is the substrate: edge supports (routing reliability
per outgoing connection) and context-indexed action supports (transform value per
context bit). Both persist across sessions and are the basis of carryover
experiments. In \textit{visible} mode, the context bit is given directly to nodes.
In \textit{latent} mode, nodes infer context from packet flow statistics without
an explicit label. Latent inference is slower to commit but accumulates substrate
entries without context binding, making them transfer-safe by construction.

When a node accumulates ATP surplus while failing to route productively, it can
bud a new node or edge. The surplus gate means growth fires when the node is
metabolically stressed --- exactly where additional routing capacity is most
needed. Difficulty-correlated topology growth is a consequence of this gate, not
a separately designed objective. A recognition--anticipation--prediction loop,
added in the current architecture, lets each node compare incoming observations
against consolidated patterns, generate predictions, and modulate selection on
prediction error --- a local implementation of predictive processing
\citep{rao1999predictive} with no global consistency requirement.

\section{Interpretability}

\subsection{As Structural Capability}

REAL is natively interpretable because computation is organized as local agents
making explicit choices rather than as global function composition. Three
consequences follow.

\textbf{Selection is readable as a term-by-term comparison.} The selector exposes
the contributing scores for each candidate action: substrate support weight,
coherence history bias, recognition confidence, and ATP cost, each as a named
value. A routing decision is a record of what evidence supported what action ---
not a matrix multiply.

\textbf{Memory is inspectable as structure.} What a node has learned is encoded in
specific, named substrate entries. The question ``does Task A's substrate help or
hurt Task B?'' is answerable directly by reading which context-indexed action
supports carry over and whether they match or conflict with Task B's transform
requirements. This is categorically different from inspecting weight norms.

\textbf{Counterfactuals are mechanistically meaningful.} Many experiments are
designed as mechanism toggles: recognition enabled vs.\ disabled, fresh-session
vs.\ persistent evaluation, full vs.\ substrate-only carryover, visible vs.\
latent context mode. These isolate specific causal contributions. In a
weight-matrix system, disabling ``recognition'' requires architectural surgery that
changes the function computed; here, toggling a mechanism holds everything else
constant.

A further property: null results are diagnostically informative. When a mechanism
produces no improvement in aggregate metrics, local observability allows a more
precise question --- was the mechanism absent? Present but too late? Present but
competing with a stronger signal and losing? This precision prevents false
conclusions from aggregate accuracy alone.

\subsection{As Development Workflow}

This structural capability was the primary engineering methodology throughout
development. The recurring pattern was: inspect local state $\rightarrow$ localize
the specific failure $\rightarrow$ change that mechanism $\rightarrow$ verify. We
document three cases.

\textbf{Case 1: Downstream transform commitment under latent uncertainty.} One
task variant underperformed a closely related variant despite identical behavior at
the first latent node. A per-node probe revealed that a downstream node was
repeatedly selecting transform-specific routes while latent context confidence was
well below the promotion threshold --- in one representative trace, choosing a hard
transform at cycles where confidence was 0.289, then 0.477, then 0.609, none of
which crossed the 0.78 threshold for context promotion. The diagnosis was precise:
the node had a packet and a task but was committing to transforms before it had
sufficient contextual evidence. The intervention was targeted: a downstream action
gate that suppresses non-identity transform choices while latent context is
available but unpromoted. Task performance improved sharply in subsequent
benchmarks. What the workflow provided was not ``change the learning rate'' or
``train longer,'' but a specific mechanistic claim followed by a minimal surgical
fix.

\textbf{Case 2: Recognition absent, then present, then blunt.} An early probe
showed zero recognized route entries during warm transfer. Rather than concluding
recognition was ineffective, the trace was inspected: route patterns were
consolidated in substrate key space, but transfer-time observations did not expose
those same keys, so the recognizer had no aligned coordinate system to compare
against. After aligning the observation adapter to expose substrate-key indices,
recognized route entry count moved from zero to five --- recognition was now
genuinely active. A subsequent selector interaction probe then showed that sparse
recognition was occasionally winning close routing decisions, but tipping them
toward stale carried route families. The conclusion was not ``disable recognition''
but ``recognition needs freshness discounting.'' Three distinct diagnostic
conclusions --- absent, present but misaligned, present but blunt --- emerged from
a sequence of local observations, each narrowing the question.

\textbf{Case 3: Occupancy V1 failure diagnosed at the mechanism level.} The
initial occupancy evaluation produced an eval F1 of 0.032. The traces were
inspected rather than treating this as architectural incapacity. Three compounding
failure modes were found: feedback was suppressed during evaluation, causing ATP to
collapse as nodes spent their budgets without replenishment --- dropped packets
surged from under 1 per episode during training to over 17 per episode during
evaluation; no carryover round-trip was ever measured; and the context-indexed
support layer was dormant because no context signal was injected. Addressing the
first flaw alone --- enabling feedback during evaluation --- raised eval F1 from
0.032 to 0.748, a 23$\times$ improvement, establishing that the failure was
entirely in the evaluation harness.

\section{Experimental Results}

\subsection{CVT-1: Transfer, Latent Context, and Sample Efficiency}

The primary controlled benchmark uses context-dependent bit-transform routing
tasks (CVT-1). Three variants (A, B, C) form a structured transfer landscape where
some prior substrate provides a useful prior and some injects context poison.

\begin{table}[ht]
  \centering
  \small
  \begin{tabular}{lcc}
    \toprule
    Condition & Exact / 18 & Bit acc. \\
    \midrule
    Task B cold start          & 3.9  & 0.477 \\
    A$\to$B substrate-only     & 7.3  & 0.586 \\
    A$\to$B full carryover     & 10.6 & 0.704 \\
    A$\to$B latent carryover   & 7.0  & ---   \\
    A$\to$B visible carryover  & 6.2  & ---   \\
    \bottomrule
  \end{tabular}
  \caption{CVT-1 transfer results, 6-node topology, 18-packet sessions.}
  \label{tab:cvt1}
\end{table}

Both forms of carryover substantially outperform cold start, confirming the
substrate is doing real work. Latent A$\to$B transfer (7.0 exact) slightly exceeds
visible (6.2) because latent substrate entries carry no context binding and
therefore cannot introduce poison into the new task's transform mapping.

Sequential A$\to$B$\to$C transfer reaches 7.6 exact matches on Task C, equal to a
direct A$\to$C skip but via different substrate mechanisms. Neither path overwrites
prior structure: the substrate specializes through layered accumulation. An online
neural baseline comparison (same 18-example session, predict-then-update) shows an
Elman RNN requires approximately 144--162 examples to reach the criterion REAL
achieves in a single 18-packet session --- an 8--9$\times$ sample efficiency
advantage in the low-data regime.

\subsection{Morphogenesis: Emergent Difficulty-Correlated Growth}

On a 10-node topology with 36-packet sessions, morphogenesis produces outcomes
that differ systematically by baseline task performance:

\begin{table}[ht]
  \centering
  \small
  \begin{tabular}{lccc}
    \toprule
    Task & Fixed & Growth & $\Delta$ \\
    \midrule
    A (high baseline: 14.8/36) & 14.8 & 13.2 & $-$1.6 \\
    B (low baseline: 11.8/36)  & 11.8 & 19.2 & \textbf{+7.4} \\
    C (mod.\ baseline: 17.0/36) & 17.0 & 16.0 & $-$1.0 \\
    \bottomrule
  \end{tabular}
  \caption{Morphogenesis on 10-node topology. Growth helps where performance is
  lowest and disrupts where it is already high.}
  \label{tab:morpho}
\end{table}

This pattern is not designed. High-performing tasks consume ATP efficiently,
leaving little surplus to trigger growth. Struggling tasks fail more, spending ATP
on misrouted packets and opening the surplus window that enables budding. New nodes
then provide capacity precisely where the network is metabolically stressed.
Earned growth rate improved from 20\% on the 6-node topology to 100\% on the
10-node topology under task-carrying conditions, revealing a minimum topology
complexity threshold: below it, the edge space is too constrained for new nodes to
find productive specializations before the session ends.

\subsection{Occupancy: Real-World Generalization}

The occupancy series applies REAL to a 14-day room occupancy dataset: five
environmental sensors (CO$_2$, temperature, humidity, light, humidity ratio) at
15-minute intervals predict room occupancy. This is an external real-world dataset
the system was not designed for.

Following the V1 harness diagnosis in Case 3, the V3 harness aligned the
evaluation to REAL's native operating mode: sessions as first-class units, online
multi-sensor context encoding (CO$_2$ and light combined into a 4-code context
space), full multi-hop routing, and explicit fresh-session vs.\ persistent
evaluation protocols. Results across three seeds:

\begin{table}[ht]
  \centering
  \small
  \begin{tabular}{lccc}
    \toprule
    Condition & Acc. & F1 & Delivery \\
    \midrule
    MLP baseline (full data)        & 0.985 & 0.963 & ---    \\
    REAL V3 warm (persistent)       & 0.973 & \textbf{0.952} & 0.969 \\
    REAL V3 cold (fresh-session)    & 0.968 & 0.919 & 0.977  \\
    Efficiency ratio (cold/warm del.) & --- & ---   & \textbf{0.991} \\
    \bottomrule
  \end{tabular}
  \caption{V3 occupancy results, 3-seed mean. MLP trained on full 1,344-example
  split; REAL trained on 926 sessions.}
  \label{tab:occupancy}
\end{table}

Warm F1 of 0.952 closes the gap to the MLP to 0.011. The remaining difference is
a recall gap on the minority class attributable to the absence of class-weighted
feedback; both systems achieve near-identical precision ($\sim$0.98). The
efficiency ratio of 0.991 means a fresh system loaded with the trained substrate
routes at 99.1\% of continuous-system delivery --- from the first session, without
warm-up --- confirming that the substrate is the actual knowledge carrier.

\section{Discussion}

\subsection{What Building the System Revealed}

ALIFE's framing of artificial life as experimental philosophy is directly
applicable here. Several findings emerged from operating REAL that prior
formalization did not anticipate.

The most significant is that the right evaluation frame for this kind of system is
not obvious and cannot be derived from the theory. The occupancy V1 result ---
0.032 eval F1 --- appeared to be a system failure. Node-level traces showed
instead a harness failure: the evaluation was actively starving the network by
suppressing the feedback nodes require to maintain their ATP. When deprived of
feedback, the system did exactly what an allostatic system should do: it conserved
remaining budget and stopped routing. This was not incorrect behavior. It was the
architecture working as designed under conditions the harness had not anticipated.
The failure was in the measurement, not in the system, and only local observability
made that distinction possible.

A second discovery concerns substrate portability. The original research question
was ``how does REAL compare to an MLP on classification accuracy?'' The more
fundamental finding is that a fresh system loaded with trained substrate performs
within 0.85\% of a continuously running system from the very first session. The
routing substrate is not a snapshot of runtime state; it is a compressed record of
which routing patterns have earned feedback, indexed by context, that transfers
like a learned reflex rather than like a weight matrix. This property was not
predicted by the theory; it emerged from measuring the right thing.

A third discovery is the difficulty-correlated morphogenesis pattern. The ATP
surplus gate was designed to prevent runaway growth. That it also produces
self-allocation of structural resources tracking the network's own deficits ---
without being told to --- was not predicted. It is the metabolic constraint
producing behavior that resembles a biological property (allocate growth where
stressed) through a mechanism that is entirely local and substrate-level.

These three discoveries share a structure: they were accessible only through
operating the system and reading its local state. None appear in the formal
specification. This is the experimental philosophy claim: the built artifact makes
falsifiable claims the formal system does not, and some of those claims turn out to
be correct in ways that revise the theory.

\subsection{On the Two Roles of Interpretability}

The capability claim and the workflow claim are separable and should not be
conflated. The capability claim is structural: REAL decisions are readable because
they are made by agents with explicit local state. When a node selects an action,
there is a node, a state it observed, a set of candidate actions with named scores,
and a selected action. All are first-class objects.

The workflow claim is empirical: this readability was used. Across three documented
cases, the methodology was inspect $\to$ localize $\to$ change $\to$ verify, and
each step was possible because of the structural capability. A transformer with
richer logging would not produce the claim ``a downstream node is committing to
transforms before context confidence is sufficient'' --- because there is no
downstream node, there is no context confidence state, and there is no action
vocabulary from which a premature commitment was made.

\subsection{Limits}

Morphogenesis requires task and context metadata; load-based growth in the absence
of task structure remains outside the current design. Prediction's influence on
routing decisions is currently observable but weak in the tested session lengths.
The occupancy recall gap on the minority class would likely close with
class-weighted feedback. The sample-efficiency advantage at small scale has not
been tested at the scales where gradient-based methods dominate.

\section{Related Work}

\textbf{Biologically-plausible learning.} REAL connects to local learning rules
that avoid global credit assignment, including contrastive Hebbian learning
\citep{movellan1991contrastive} and equilibrium propagation
\citep{scellier2017equilibrium}. It differs in replacing the optimization objective
with allostatic coherence rather than approximating backprop with a local rule.

\textbf{Stigmergy and swarm systems.} The substrate-as-memory architecture has
structural parallels with ant colony optimization \citep{dorigo1996antsystem} and
stigmergic multi-agent systems \citep{theraulaz1999stigmergy}. REAL differs in
orienting toward generalization and cross-task transfer rather than combinatorial
optimization.

\textbf{Reservoir computing.} Liquid state machines and echo state networks also
use substrate structure for computation \citep{maass2002realtime,jaeger2001echo}.
Their substrates are random or fixed rather than learned through metabolic
feedback, transfer is not a design objective, and local decisions are no more
interpretable than in standard networks.

\textbf{Interpretable machine learning.} Post-hoc methods \citep{lundberg2017shap,
ribeiro2016lime, sundararajan2017integrated} approximate explanations of trained
models and are correlational by construction. REAL's interpretability is causal:
the explanation of a routing decision is the decision record itself, not an
approximation derived from perturbing the input space.

\section{Conclusion}

REAL Neural Substrate demonstrates that a network of local allostatic agents,
without global gradient descent, can develop cross-task transfer, latent context
inference, difficulty-correlated topology growth, and near-MLP performance on a
real-world dataset. The biological principles it operationalizes --- metabolic
constraint, allostatic coherence, stigmergic substrate --- are sufficient to
produce competitive behavior without a global objective.

The system is natively interpretable in two distinct senses. As a structural
capability, REAL decisions are readable because they are made by agents with
explicit local state; memory is inspectable because it is encoded in named
substrate entries; and counterfactual analysis is mechanistically meaningful
because mechanisms can be toggled without redesigning the computation. As a
development workflow, this readability enabled an approach --- inspect, localize,
change, verify --- that was productive across the documented cases and would not
have been available in a gradient-trained system.

Three experimental discoveries --- that the occupancy failure was a harness failure
readable from metabolic traces, that substrate portability is the more fundamental
capability, and that difficulty-correlated growth is an emergent property of the
metabolic gate --- were all accessible only through building and operating the
system. They constitute the experimental-philosophy contribution: the built
artifact makes claims the formal specification does not, and some of those claims
revise the theory.

The interpretability is not a feature. It is what the architecture makes possible.

\section*{Acknowledgements}

The author thanks the Codex coding agent for implementation support throughout
Phase 8 development, and Gemini for the generative question that initiated the
neural substrate direction.

\footnotesize
\bibliographystyle{apalike}
\bibliography{references}

\end{document}
